\section{Трекинг модельного кольца}
\label{app:ch3track}

Кольцо раздела~\ref{sec:es-num} собрано из восьми одинаковых ячеек. Опорная дуга --- четверть окружности радиуса $\rho=20\,\mathrm{м}$. Прямой промежуток ячейки равен $8\,\mathrm{м}$ и делится квадруполем пополам. Знак тонкой линзы чередуется, модуль силы равен $0{,}02\,\mathrm{м}^{-1}$. Секступоль, когда он включён, стоит на фокусирующих плоскостях и даёт толчок $\Delta x'=(S/2)x^2(p_s/p)$. Цилиндрическое поле внутри дуги имеет потенциал~\eqref{eq:es-phi}; на границе дуги продольный импульс меняется так, чтобы полная энергия~\eqref{eq:es-energy} и радиальный импульс сохранились. Магнитный вариант заменяет дугу однородным полем $B=p_s/(q\rho)$ и масштабирует квадруполь как $1/p$, а не как $1/(pv)$. Опорная частица --- протон с кинетической энергией $233\,\mathrm{МэВ}$.

Замкнутая орбита при заданном $\delta$ ищется в плоскости первого квадруполя. Коэффициенты $\alpha_0$, $\alpha_1$ берутся из длины этой орбиты, $\eta_0$ и $\eta_1$ --- из времени. Подгонка кубическим полиномом по набору $|\delta|\le 10^{-3}$ оставляет остаток меньше $10^{-12}$. Хроматичность есть наклон тюна замкнутой орбиты по $\delta$. Бетатронный наклон есть среднее $(T-T_s)/T_s$ и $(C-C_s)/C_s$ по двадцати четырём начальным фазам линейного эллипса Твисса, делённое на действие $J$. При $J=10^{-5}\,\mathrm{м}$ и при $J=4\times 10^{-5}\,\mathrm{м}$ магнитные наклоны времени и пути совпадают до $2\times 10^{-7}\,\mathrm{м}^{-1}$. Расхождение этих двух наклонов в прежней точке $J=10^{-8}\,\mathrm{м}$ было сравнимо с допуском интегратора. Слежение одной частицы на $40$--$80$ оборотах постоянную часть от бетатронного колебания периода не отделяет: амплитуда колебания больше среднего в сотни раз, и среднее по такому отрезку зависит от числа оборотов. В таблицу входит фазовое среднее.

Определитель матрицы $(x,x')$ на опорной орбите отличается от единицы меньше чем на $2\times 10^{-6}$.

\begin{table}[ht]
\centering
\caption{Наклоны периода по действию. Предсказание есть $-2\pi\xi/C_s$.}
\label{tab:es-scan}
\begin{tabular}{llrrrrr}
\hline
кольцо & $S$ & $\xi$ & время & путь & $-2\pi\xi/C_s$ \\
\hline
магнит & 0 & $-0{,}20683$ & $8{,}36\times 10^{-3}$ & $8{,}36\times 10^{-3}$ & $6{,}85\times 10^{-3}$ \\
магнит & 0{,}002 & $+0{,}23242$ & $-7{,}28\times 10^{-3}$ & $-7{,}28\times 10^{-3}$ & $-7{,}70\times 10^{-3}$ \\
магнит & 0{,}004 & $+0{,}67071$ & $-2{,}29\times 10^{-2}$ & $-2{,}31\times 10^{-2}$ & $-2{,}22\times 10^{-2}$ \\
цилиндр & 0 & $-0{,}52236$ & $+1{,}16\times 10^{-2}$ & $-7{,}57\times 10^{-3}$ & $+1{,}73\times 10^{-2}$ \\
цилиндр & 0{,}002 & $-0{,}09641$ & $-5{,}76\times 10^{-3}$ & $-1{,}82\times 10^{-2}$ & $+3{,}19\times 10^{-3}$ \\
цилиндр & 0{,}0025 & $+0{,}01244$ & $-1{,}01\times 10^{-2}$ & $-2{,}08\times 10^{-2}$ & $-4{,}12\times 10^{-4}$ \\
цилиндр & 0{,}004 & $+0{,}33193$ & $-2{,}32\times 10^{-2}$ & $-2{,}88\times 10^{-2}$ & $-1{,}10\times 10^{-2}$ \\
цилиндр & 0{,}008 & $+1{,}18710$ & $-5{,}79\times 10^{-2}$ & $-4{,}99\times 10^{-2}$ & $-3{,}93\times 10^{-2}$ \\
цилиндр & 0{,}012 & $+2{,}04257$ & $-9{,}27\times 10^{-2}$ & $-7{,}11\times 10^{-2}$ & $-6{,}77\times 10^{-2}$ \\
\hline
\end{tabular}
\end{table}

Наклоны даны в м$^{-1}$. Прямая~\eqref{eq:es-scan} по магнитному скану имеет $\kappa=1{,}076$ и $b=1{,}00\times 10^{-3}\,\mathrm{м}^{-1}$ для времени; остаток точек меньше $6\times 10^{-6}\,\mathrm{м}^{-1}$. По цилиндру для времени $\kappa_T=1{,}227$, $b_T=-9{,}65\times 10^{-3}\,\mathrm{м}^{-1}$; для пути $\kappa_C=0{,}747$, $b_C=-2{,}05\times 10^{-2}\,\mathrm{м}^{-1}$. Точка $S=0{,}0025$ в подгонку времени входит и на прямой лежит.

Угловой интеграл $(1/(2C_s))\int\gamma_x\,\mathrm{d}s$ на магнитной оптике, совпадающей с трекингом по $\nu_x=0{,}286$ и $\beta_x=32{,}4\,\mathrm{м}$ в плоскости квадруполя, равен $2{,}41\times 10^{-2}\,\mathrm{м}^{-1}$.

Квадратичный коэффициент цилиндра при $S=0$ равен $0{,}75898$, магнитная формула даёт $0{,}87339$. Моменты $\langle(D/\rho)^2\rangle=0{,}61516$ и $\langle d_2\rangle$, экстраполированные к $\delta=0$ с орбит $|\delta|=(2\ldots 8)\times 10^{-4}$, подставленные в~\eqref{eq:es-eta1}, оставляют остаток $1\times 10^{-4}$. При $S=0{,}004$ измеренное $\eta_1^{\mathrm{es}}=0{,}31909$, магнитная формула даёт $0{,}60544$, и~\eqref{eq:es-eta1} на тех же моментах без экстраполяции $\delta^4$ отличается на $1\times 10^{-3}$.
